\chapter{Compact Operators and the Spectral Theorem}
\label{ch:b3:spectral}

Diagonalization is finite-dimensional linear algebra's crown
jewel: a symmetric matrix has an orthonormal basis of
eigenvectors. In infinite dimension this fails for bounded
self-adjoint operators in general --- multiplication by $x$ on
$L^2(\intcc01)$ has \emph{no} eigenvalues at all
(\cref{exo:b3:spectral:6}) --- but it survives, in perfect
form, for the operators that are \emph{almost finite
dimensional}: the compact ones. The spectral theorem for
compact self-adjoint operators is the single most used theorem
of applied functional analysis: it diagonalizes integral
equations, drives the Fredholm alternative, and (weekend
problem) solves the vibrating string, producing the sine basis
of Fourier analysis from pure operator theory --- with Euler's
$\zeta(2) = \frac{\pi^2}6$ falling out of a trace formula as a
parting gift. Throughout, $H$ is a Hilbert space over $\C$ (or
$\R$; statements adapt), and operators are bounded.

\section{Compact operators}

\begin{definition}\label{def:b3:spectral:compact}
$T \in \mathcal L(E, F)$ ($E, F$ Banach) is
\emph{compact}\index{compact operator} if the image $T(B)$ of
the unit ball is relatively compact in $F$ --- equivalently,
every bounded sequence $(x_n)$ has a subsequence with
$(Tx_{n_k})$ convergent. Finite-rank operators are compact
(bounded sets in finite dimension); the identity of an
infinite-dimensional space never is (Riesz's theorem, Year 2).
\end{definition}

\begin{proposition}\label{prop:b3:spectral:ideal}
The compact operators $\mathcal K(E, F)$ form a closed subspace
of $\mathcal L(E, F)$, and a two-sided ideal: $S$ compact
$\Rightarrow$ $AS$ and $SB$ compact for bounded $A, B$.
Moreover, in a Hilbert space, every compact operator is a
norm-limit of finite-rank operators.
\end{proposition}

\begin{proof}
Subspace: clear from the sequence characterization. Ideal:
bounded maps send convergent sequences to convergent ones and
bounded to bounded. Closedness: let $T_n \to T$ with $T_n$
compact, and $(x_k)$ bounded by $1$; a diagonal extraction
makes $(T_nx_{k_j})_j$ convergent for every $n$; then
$(Tx_{k_j})$ is Cauchy, since
\[
\norm{Tx_{k_j} - Tx_{k_l}} \leq 2\vertiii{T - T_n} +
\norm{T_nx_{k_j} - T_nx_{k_l}} ,
\]
choosing first $n$ then the indices. Approximation in Hilbert
spaces: let $T$ be compact, $K = \overline{T(B)}$ compact;
given $\varepsilon$, cover $K$ by finitely many balls
$B(y_i, \varepsilon)$ and let $P$ be the orthogonal projection
onto $V = \operatorname{Vect}(y_1, \dots, y_m)$ (closed:
finite-dimensional). Then $PT$ has finite rank, and for
$\norm x \leq 1$: picking $y_i$ with $\norm{Tx - y_i} <
\varepsilon$,
\[
\norm{Tx - PTx} \leq \norm{Tx - y_i} + \norm{P(y_i - Tx)}
\leq 2\varepsilon
\]
($y_i = Py_i$; $\vertiii P \leq 1$): $\vertiii{T - PT} \leq
2\varepsilon$.
\end{proof}

\begin{example}\label{ex:b3:spectral:examples}
(a) Diagonal operators on $\ell^2$: $T(x_n) = (\lambda_nx_n)$
is compact iff $\lambda_n \to 0$
(\cref{exo:b3:spectral:2}).
(b) Kernel operators on $\mathcal C(\intcc01)$: compact by
Ascoli (\cref{exo:b3:complete:7}).
(c) \emph{Hilbert--Schmidt operators}\index{Hilbert--Schmidt
operator}: for $k \in L^2(\intcc01^2)$,
\[
(T_kf)(x) = \int_0^1k(x, y)\,f(y)\,\dd y
\]
defines a compact operator on $L^2(\intcc01)$ with
$\vertiii{T_k} \leq \norm k_{L^2}$
(\cref{exo:b3:spectral:4}: truncating the basis expansion of
$k$ exhibits $T_k$ as a limit of finite-rank operators).
\end{example}

\section{Self-adjoint operators}

\begin{definition}\label{def:b3:spectral:selfadjoint}
$T \in \mathcal L(H)$ is \emph{self-adjoint}\index{self-adjoint
operator} if $T = T^*$ (\cref{exo:b3:hilbert:8}), i.e.\
$\langle Tx, y\rangle = \langle x, Ty\rangle$ for all $x, y$.
Then $\langle x, Tx\rangle \in \R$ for every $x$ (equal to its
conjugate).
\end{definition}

\begin{proposition}\label{prop:b3:spectral:sanorm}
For self-adjoint $T$:\index{numerical radius}
\[
\vertiii T = \sup_{\norm x \leq 1}\ \abs{\langle x,
Tx\rangle} .
\]
Eigenvalues of $T$ are real, and eigenvectors for distinct
eigenvalues are orthogonal.
\end{proposition}

\begin{proof}
Let $M$ be the supremum; $M \leq \vertiii T$ by
Cauchy--Schwarz. Conversely, the polarization-type identity
\[
\langle x{+}y, T(x{+}y)\rangle - \langle x{-}y,
T(x{-}y)\rangle = 4\operatorname{Re}\langle y, Tx\rangle
\]
(expand; the cross terms $\langle y, Tx\rangle + \langle x,
Ty\rangle = 2\operatorname{Re}\langle y, Tx\rangle$ by
self-adjointness) gives, with the parallelogram law,
\[
4\operatorname{Re}\langle y, Tx\rangle \leq
M\bigl(\norm{x{+}y}^2 + \norm{x{-}y}^2\bigr) = 2M\bigl(\norm
x^2 + \norm y^2\bigr).
\]
For $\norm x = 1$ with $Tx \neq 0$, take $y = Tx/\norm{Tx}$:
$4\norm{Tx} \leq 4M$. So $\vertiii T \leq M$. Eigenvalues:
$Tx = \lambda x$, $x \ne 0$ gives $\lambda\norm x^2 = \langle
x, Tx\rangle \in \R$. Orthogonality: $\lambda\langle x,
y\rangle = \langle Tx, y\rangle = \langle x, Ty\rangle =
\mu\langle x, y\rangle$ with $\lambda \neq \mu$ real.
\end{proof}

\section{The spectral theorem}

\begin{lemma}[Existence of an extreme eigenvalue]
\label{lem:b3:spectral:existence}
Let $T \neq 0$ be compact and self-adjoint. Then $\vertiii T$
or $-\vertiii T$ is an eigenvalue of $T$.
\end{lemma}

\begin{proof}
By \cref{prop:b3:spectral:sanorm}, pick unit vectors $x_n$
with $\langle x_n, Tx_n\rangle \to \mu$, where $\abs\mu =
\vertiii T > 0$ (pass to a subsequence to fix the sign). Then
\[
\norm{Tx_n - \mu x_n}^2
= \norm{Tx_n}^2 - 2\mu\langle x_n, Tx_n\rangle + \mu^2
\leq \vertiii T^2 - 2\mu\langle x_n, Tx_n\rangle + \mu^2
\longrightarrow 2\mu^2 - 2\mu\cdot\mu = 0 .
\]
By compactness, a subsequence $Tx_{n_k} \to y$; then $\mu
x_{n_k} = Tx_{n_k} - (Tx_{n_k} - \mu x_{n_k}) \to y$, so
$x_{n_k} \to x = y/\mu$, a unit vector, and continuity gives
$Tx = \mu x$.
\end{proof}

\begin{theorem}[Spectral theorem for compact self-adjoint
operators]\label{thm:b3:spectral:spectral}
Let $T$ be a compact self-adjoint operator on a Hilbert space
$H$.\index{spectral theorem (compact self-adjoint)}
\begin{enumerate}
  \item $H$ admits an orthonormal system $(e_n)_{n \in N}$
        ($N$ finite or countable) of eigenvectors of $T$, with
        real nonzero eigenvalues $(\lambda_n)$, such that
        \[
        Tx = \sum_{n\in N}\lambda_n\,\langle e_n, x\rangle\,
        e_n \qquad (x \in H),
        \]
        and $H = \ker T \,\oplus^\perp\,
        \overline{\operatorname{Vect}}(e_n : n \in N)$.
  \item If $N$ is infinite, $\lambda_n \to 0$; for each
        $\delta > 0$ only finitely many $n$ have
        $\abs{\lambda_n} \geq \delta$, and each eigenspace
        $\ker(T - \lambda)$, $\lambda \neq 0$, is
        finite-dimensional.
  \item Completing $(e_n)$ by an orthonormal basis of $\ker T$
        yields, when $H$ is separable, an orthonormal basis of
        $H$ made of eigenvectors: $T$ is diagonalized.
\end{enumerate}
\end{theorem}

\begin{proof}
(2) first. If infinitely many orthonormal eigenvectors $x_k$
had $\abs{\lambda_{(k)}} \geq \delta$: $\norm{Tx_k - Tx_l}^2 =
\lambda_{(k)}^2 + \lambda_{(l)}^2 \geq 2\delta^2$
(orthogonality, Pythagoras): no convergent subsequence of
$(Tx_k)$, contradicting compactness of $T$ on the bounded
$(x_k)$. This bounds by a finite number, for each $\delta$,
the total multiplicity of eigenvalues outside
$\intoo{-\delta}\delta$; countability and $\lambda_n \to 0$
follow.

(1) Let $H_0$ be the closed span of \emph{all} eigenvectors
with nonzero eigenvalues, organized (by (2) and
Gram--Schmidt within each finite-dimensional eigenspace,
orthogonality across eigenspaces from
\cref{prop:b3:spectral:sanorm}) into an orthonormal system
$(e_n)$ with eigenvalues $\lambda_n \neq 0$. $T$ maps $H_0$
into $H_0$, and also $H_0^\perp$ into $H_0^\perp$: for $y
\perp H_0$ and $e$ an eigenvector, $\langle e, Ty\rangle =
\langle Te, y\rangle = \lambda\langle e, y\rangle = 0$. The
restriction $T' = T\restriction_{H_0^\perp}$ is compact
self-adjoint on the Hilbert space $H_0^\perp$; if $T' \neq
0$, \cref{lem:b3:spectral:existence} produces an eigenvector
of $T$ with nonzero eigenvalue inside $H_0^\perp$ ---
impossible, such vectors live in $H_0$. So $T' = 0$:
$H_0^\perp \subseteq \ker T$. Conversely $\ker T \perp$ every
$e_n$ ($\langle e_n, z\rangle = \frac1{\lambda_n}\langle
Te_n, z\rangle = \frac1{\lambda_n}\langle e_n, Tz\rangle =
0$): $\ker T \subseteq H_0^\perp$, whence $\ker T =
H_0^\perp$ and the orthogonal decomposition. The expansion:
for $x = z + \sum_nc_ne_n$ ($z \in \ker T$, $c_n = \langle
e_n, x\rangle$; \cref{thm:b3:hilbert:parseval}(1) on $H_0$),
continuity of $T$ gives $Tx = \sum_nc_n\lambda_ne_n$.

(3) $\ker T$, a closed subspace of a separable space, is
separable: it has an orthonormal basis
(\cref{prop:b3:hilbert:gramschmidt}); the union is an
orthonormal basis of $H$ by the decomposition in (1).
\end{proof}

\begin{theorem}[Fredholm alternative]
\label{thm:b3:spectral:fredholm}
Let $T$ be compact self-adjoint and $\lambda \in \R\setminus
\{0\}$.\index{Fredholm alternative}
\begin{enumerate}
  \item If $\lambda$ is not an eigenvalue, then $T - \lambda I$
        is bijective with bounded inverse: for every $f$, the
        equation $Tx - \lambda x = f$ has exactly one
        solution, depending continuously on $f$.
  \item If $\lambda$ is an eigenvalue, $Tx - \lambda x = f$ is
        solvable iff $f \perp \ker(T - \lambda I)$, and the
        solution is unique up to that (finite-dimensional)
        kernel.
\end{enumerate}
\end{theorem}

\begin{proof}
Decompose $x = z + \sum c_ne_n$ and $f = w + \sum d_ne_n$
along \cref{thm:b3:spectral:spectral} ($z, w \in \ker T$). The
equation reads
\[
-\lambda z = w, \qquad (\lambda_n - \lambda)\,c_n = d_n\
(n \in N).
\]
(1) $\lambda \notin \{\lambda_n\}\cup\{0\}$: by (2) of the
spectral theorem, $\inf_n\abs{\lambda_n - \lambda} = \delta >
0$ (eigenvalues accumulate only at $0 \neq \lambda$). Solve:
$z = -w/\lambda$, $c_n = d_n/(\lambda_n - \lambda)$, with
$\sum\abs{c_n}^2 \leq \delta^{-2}\sum\abs{d_n}^2$: a unique
solution with $\norm x \leq C\norm f$.
(2) $\lambda = \lambda_{n}$ for $n$ in a finite set $F$:
solvability of $(\lambda_n - \lambda)c_n = d_n$ for $n \in F$
requires $d_n = 0$, i.e.\ $f \perp e_n$ ($n \in F$), i.e.\ $f
\perp \ker(T - \lambda I)$; the $c_n$, $n \in F$, are then
free.
\end{proof}

\begin{example}\label{ex:b3:spectral:minxy}
On $L^2(\intcc01)$, let $Tf(x) = \int_0^1\min(x, y)f(y)\dd y$:
a Hilbert--Schmidt operator with real symmetric kernel:
compact and self-adjoint. Solving $Tf = \lambda f$: the
relation $\bigl(Tf\bigr)(x) = \int_0^xyf(y)\dd y +
x\int_x^1f(y)\dd y$ shows $u = Tf$ satisfies $u'' = -f$ (two
differentiations, legitimate for continuous $f$, and $Tf$ is
continuous for $f \in L^2$: dominated convergence), with $u(0)
= 0$ and $u'(1) = 0$. So eigenfunctions solve $\lambda u'' =
-u$, $u(0) = 0$, $u'(1) = 0$:
\[
u_n(x) = \sin\Bigl(\bigl(n + \tfrac12\bigr)\pi x\Bigr),
\qquad
\lambda_n = \frac{1}{\bigl(n + \frac12\bigr)^2\pi^2}
\quad (n \geq 0),
\]
and the spectral theorem asserts --- with no Fourier theory ---
that these sines form an orthonormal basis of $L^2(\intcc01)$
after normalization (the kernel of $T$ is $0$: $Tf = 0$
forces, by the two differentiations, $f = 0$ a.e.). The
weekend problem runs the same circle of ideas for the
vibrating string and extracts $\zeta(2)$ from the trace.
\end{example}

\begin{method}\label{met:b3:spectral:method}
Given an integral or differential equation: (1) recast it as
$(I - \lambda K)u = f$ or $Ku = \lambda u$ with $K$ an
integral operator; (2) verify $K$ compact (Hilbert--Schmidt
kernel, or Ascoli) and, if possible, self-adjoint (symmetric
real kernel); (3) diagonalize with the spectral theorem or
invoke the Fredholm alternative for solvability; (4) read off
existence, uniqueness, stability, and series formulas for
solutions in the eigenbasis. Differential operators are
unbounded, but their \emph{inverses} (Green operators) are
compact: always invert first.
\end{method}

\section{Exercises}

\begin{exercise}[$\star$]\label{exo:b3:spectral:1}
(a) Show that a bounded operator with finite-dimensional range
is compact.
(b) Show that the identity of a normed space is compact iff
the dimension is finite (Riesz, Year 2). Deduce that a compact
operator on an infinite-dimensional space is never invertible
with bounded inverse.
\end{exercise}

\begin{exercise}[$\star$]\label{exo:b3:spectral:2}
Let $T(x_1, x_2, \dots) = (\lambda_1x_1, \lambda_2x_2, \dots)$
on $\ell^2$, with $(\lambda_n)$ bounded.
(a) Show $\vertiii T = \sup\abs{\lambda_n}$.
(b) Show that $T$ is compact iff $\lambda_n \to 0$.
\emph{(For $\Leftarrow$, truncate; for $\Rightarrow$, test on
$(e_n)$.)}
(c) When is $T$ self-adjoint? Verify the spectral theorem by
inspection in that case.
\end{exercise}

\begin{exercise}[$\star\star$]\label{exo:b3:spectral:3}
Give the details of the ideal property
(\cref{prop:b3:spectral:ideal}): if $S$ is compact and $A, B$
bounded, then $ASB$ is compact. Deduce that if $ST = TS = I$
for some bounded $S$, and $\dim H = \infty$, then $T$ is not
compact --- and reconcile with \cref{exo:b3:spectral:1}(b).
\end{exercise}

\begin{exercise}[$\star\star$]\label{exo:b3:spectral:4}
(Hilbert--Schmidt) Let $k \in L^2(\intcc01^2)$ and $(e_n)$ a
Hilbert basis of $L^2(\intcc01)$.
(a) Show that $\vertiii{T_k} \leq \norm k_{L^2}$
\emph{(Cauchy--Schwarz in the $y$-variable, then Tonelli)}.
(b) Expand $k(x,y) = \sum_{m,n}c_{mn}e_m(x)\overline{e_n(y)}$
in $L^2$ of the square (justify that the products form a
Hilbert basis there), and show that truncating the sum gives
finite-rank operators converging to $T_k$ in operator norm:
$T_k$ is compact.
\end{exercise}

\begin{exercise}[$\star\star$]\label{exo:b3:spectral:5}
Let $T$ be self-adjoint with $\langle x, Tx\rangle \geq 0$ for
all $x$ (\emph{positive} operator).
(a) Show that eigenvalues are $\geq 0$ and that $\vertiii T =
\sup_{\norm x\leq1}\langle x, Tx\rangle$.
(b) Prove the generalized Cauchy--Schwarz inequality
$\abs{\langle x, Ty\rangle}^2 \leq \langle x, Tx\rangle\langle
y, Ty\rangle$.
\end{exercise}

\begin{exercise}[$\star\star$]\label{exo:b3:spectral:6}
On $L^2(\intcc01)$, let $(Mf)(x) = x\,f(x)$.
(a) Show that $M$ is bounded, self-adjoint, with $\vertiii M =
1$, but has \emph{no} eigenvalues.
(b) Show $M$ is not compact \emph{(exhibit a bounded sequence
whose image has no convergent subsequence, e.g.\ normalized
indicators of shrinking intervals near $1$ --- or invoke the
spectral theorem)}.
(c) Where does \cref{lem:b3:spectral:existence}'s proof break
for $M$?
\end{exercise}

\begin{exercise}[$\star\star$]\label{exo:b3:spectral:7}
(Volterra) On $L^2(\intcc01)$, let $Vf(x) = \int_0^xf(y)\dd
y$.
(a) Show $V$ is compact (Hilbert--Schmidt with kernel $\mathbf
1_{y < x}$) but not self-adjoint; compute $V^*$.
(b) Show that $V$ has no nonzero eigenvalue. \emph{(From $Vf =
\lambda f$: $f$ has a continuous representative, then is
$\mathcal C^1$, and solves $\lambda f' = f$, $f(0) = 0$.)}
(c) Conclude that compactness alone yields no eigenvectors:
self-adjointness in \cref{thm:b3:spectral:spectral} is
essential.
\end{exercise}

\begin{exercise}[$\star\star\star$]\label{exo:b3:spectral:8}
(Courant--Fischer) Let $T$ be compact, self-adjoint,
\emph{positive}, with eigenvalues $\mu_1 \geq \mu_2 \geq \dots
> 0$ (repeated by multiplicity, eigenvectors $e_1, e_2,
\dots$). Show:
\[
\mu_{k} = \max_{\substack{V \subseteq H \\ \dim V = k}}\
\min_{\substack{x \in V\\ \norm x = 1}}\ \langle x, Tx\rangle
= \min_{\substack{W \subseteq H\\ \operatorname{codim}W =
k-1}}\ \max_{\substack{x\in W\\ \norm x = 1}}\ \langle x,
Tx\rangle .
\]
\emph{(Test $V = \operatorname{Vect}(e_1,\dots,e_k)$; for the
upper bound intersect any $V$ with
$\operatorname{Vect}(e_k, e_{k+1}, \dots)$-type spaces:
dimension counting forces a nonzero intersection.)} Deduce
that eigenvalues depend monotonically on $T$ ($T \leq S
\Rightarrow \mu_k(T) \leq \mu_k(S)$).
\end{exercise}

\begin{exercise}[$\star\star$]\label{exo:b3:spectral:9}
Using \cref{thm:b3:spectral:fredholm} for $Tf(x) =
\int_0^1\min(x,y)f(y)\dd y$ (\cref{ex:b3:spectral:minxy}):
for which $\lambda \in \R$ does the integral equation
\[
f(x) - \lambda\int_0^1\min(x,y)\,f(y)\,\dd y = g(x)
\]
have a unique solution $f \in L^2$ for every $g \in L^2$? What
happens at the exceptional values?
\end{exercise}

\begin{exercise}[$\star\star\star$]\label{exo:b3:spectral:10}
Let $S$ be the shift on $\ell^2$
(\cref{exo:b3:banach:1}).
(a) Show that $S$ has no eigenvalues, while every $\lambda$
with $\abs\lambda < 1$ is an eigenvalue of $S^*$ (find the
eigenvectors explicitly, geometric sequences).
(b) Neither $S$ nor $S^*$ is compact: verify via
\cref{exo:b3:spectral:2}-style testing on $(e_n)$.
(c) Comment: for non-self-adjoint, non-compact operators, the
eigenvalue landscape can be anything from empty to a full
disc --- the notion that survives is the \emph{spectrum},
studied in a later course.
\end{exercise}

\section{Problem: the vibrating string and
\texorpdfstring{$\zeta(2)$}{zeta(2)}}

\begin{problem}[{Weekend problem --- Green's operator,
the sine basis, and a trace formula}]\label{pb:b3:spectral:1}
We solve the eigenvalue problem of the vibrating string with
fixed ends --- $-u'' = \nu u$, $u(0) = u(1) = 0$ --- by
operator theory, obtain the sine orthonormal basis without any
Fourier computation, and evaluate $\zeta(2)$ by comparing two
expressions for the \emph{trace} of the Green operator.
Define, on $L^2(\intcc01)$,
\[
(Gf)(x) = \int_0^1 g(x,y)\,f(y)\,\dd y,
\qquad
g(x, y) = \min(x,y)\,\bigl(1 - \max(x,y)\bigr).
\]

\medskip
\noindent\textbf{Part I --- The Green operator.}
\begin{enumerate}
  \item Show that $g$ is continuous, symmetric, with $0 \leq g
        \leq \frac14$, and that $G$ is compact and
        self-adjoint (\cref{ex:b3:spectral:examples}(c)).
  \item For $f$ continuous, show that $u = Gf$ is $\mathcal
        C^2$ with
        \[
        -u'' = f, \qquad u(0) = u(1) = 0
        \]
        \emph{(write $u(x) = (1-x)\int_0^xyf(y)\dd y +
        x\int_x^1(1-y)f(y)\dd y$ and differentiate twice)}.
        Conversely, if $u \in \mathcal C^2$ with $u(0) = u(1)
        = 0$, then $G(-u'') = u$: $G$ inverts the string
        operator.
  \item Show $\ker G = \{0\}$ \emph{(if $Gf = 0$ with $f \in
        L^2$: test against continuous $\varphi$, transfer $G$
        by symmetry/Fubini onto $\varphi$, and use the
        fundamental lemma \cref{cor:b3:lp:fundlemma} --- or
        regularize)}, and that $G$ is a positive operator:
        $\langle f, Gf\rangle \geq 0$. \emph{(For continuous
        $f$: $\langle f, Gf\rangle = \int_0^1 (u')^2$ with $u
        = Gf$, by parts; conclude by density.)}
\end{enumerate}

\medskip
\noindent\textbf{Part II --- Diagonalization: the sine
basis.}
\begin{enumerate}[resume]
  \item Show that the eigenfunctions of $G$ with eigenvalue
        $\lambda \ne 0$ are, up to scalars, the solutions of
        $-\lambda u'' = u$, $u(0) = u(1) = 0$ \emph{(an
        eigenfunction has a continuous representative --- $Gf$
        is continuous for $f \in L^2$, why? --- hence is
        $\mathcal C^2$ by bootstrapping question 2)}.
  \item Solve the boundary value problem: the eigenvalues of
        $G$ are $\lambda_n = \frac1{n^2\pi^2}$ ($n \geq 1$),
        with normalized eigenfunctions $e_n(x) =
        \sqrt2\,\sin(n\pi x)$; check orthonormality by direct
        integration as a sanity test.
  \item Conclude from
        \cref{thm:b3:spectral:spectral} and question 3 that
        $\bigl(\sqrt2\sin(n\pi x)\bigr)_{n\geq1}$ is an
        orthonormal \emph{basis} of $L^2(\intcc01)$ --- no
        Stone--Weierstrass, no Fourier series needed. Expand
        $f(x) = x(1-x)$ in this basis and write Parseval for
        it.
\end{enumerate}

\medskip
\noindent\textbf{Part III --- The trace formula and
$\zeta(2)$.}
\begin{enumerate}[resume]
  \item Prove the two identities
        \[
        \langle e_n, Ge_n\rangle = \lambda_n
        \quad\text{and}\quad
        \sum_{n\geq1}\lambda_n = \int_0^1 g(x,x)\,\dd x .
        \]
        For the second (the \emph{trace formula}): expand
        $g(x, \cdot)$, for fixed $x$, in the basis $(e_n)$ ---
        show that the coefficients are $\lambda_ne_n(x)$, so
        that $g(x, \cdot) = \sum_n\lambda_ne_n(x)\,e_n$ in
        $L^2$. Here the sines are explicit: verify directly
        that $\sum_n\lambda_ne_n(x)e_n(y)$ converges
        \emph{uniformly} on the square (compare with $\sum
        \frac2{n^2\pi^2}$), so its sum is continuous and, having
        the same $L^2$-expansions in $y$ for each $x$, equals
        $g(x,y)$ everywhere. Set $y = x$ and integrate term by
        term.
  \item Compute $\int_0^1g(x,x)\dd x = \int_0^1x(1-x)\dd x =
        \frac16$, and conclude
        \[
        \sum_{n\geq1}\frac{1}{n^2\pi^2} = \frac16,
        \qquad\text{i.e.}\qquad
        \boxed{\ \zeta(2) = \frac{\pi^2}{6}\ } :
        \]
        Euler's sum from an operator trace.
  \item Re-derive $\zeta(2)$ a third way: apply Parseval in
        the sine basis to the constant function $\mathbf 1$,
        compute $\sum_{n \text{ odd}}\frac1{n^2}$, and
        conclude. Then compare mechanisms: in what sense is
        the trace argument of questions 7--8 ``Parseval
        applied to the whole kernel at once''?
\end{enumerate}

\medskip
\noindent\textbf{Part IV --- The string vibrates.}
\begin{enumerate}[resume]
  \item (Separation of variables, synthesized) For $f \in
        L^2$, define
        \[
        u(t, x) = \sum_{n\geq1}\;c_n\,
        \cos(n\pi t)\,\sqrt2\sin(n\pi x),
        \qquad c_n = \langle e_n, f\rangle .
        \]
        Show the series converges in $L^2(\intcc01)$ for each
        $t$, that $t\mapsto u(t, \cdot)$ is continuous into
        $L^2$, and that for $f$ in the span of finitely many
        $e_n$ it solves the wave equation $\partial_t^2u =
        \partial_x^2u$ with $u(0) = f$, $\partial_tu(0) = 0$,
        fixed ends. The eigenvalues $n^2\pi^2$ are the squared
        frequencies: the string's harmonics --- explain the
        musical interpretation of
        \cref{thm:b3:spectral:spectral} in one paragraph.
\end{enumerate}
\end{problem}
